% 旧中心商論文VI の第14--15節：中心元の有限表示とAzumaya locusの証明依存
% このファイルは papers/central_quotient_local_models.tex から入力する。
\section{全中心の明示生成元とRees退化によるAzumaya locus}
\label{sec:quotient-center-azumaya}
$R_1$を
\[
 \CC\langle X,Y,Z\rangle/
 ([X,Y]-Z^2,\,[Y,Z]-(X^2+1),\,[Z,X]-Y^2)
\]
とする。語長filtrationを$F_\bullet R_1$と書く。
旧中心商原稿の第14節では$\operatorname{gr}R_1=S_1$、
$\dim F_nR_1=\binom{n+3}{3}$を用いて、
次の非斉次中心元を明示する。

\begin{proposition}[四つの中心生成元：有限計算に依存する結果]
\label{prop:quotient-full-center}
$\Omega_1=(X^3+Y^3+Z^3)/3-
 (XYZ+XZY+YXZ+YZX+ZXY+ZYX)/9$と置き、
$g=\frac32(\Omega_1+X)=Y^3-ZYX+X$とする。
以下の語を定義する：
\[
 U=Z^3(g-Z^3)=Y^3(g-Y^3),
\]
\[
\begin{aligned}
 V={}&YZ^2Y^3-ZYZ^2YX-ZYZ^4
   +2Z^2Y^2ZX+Z^2YZ^3-3Z^3YZ^2\\
  &+YZ^2X+ZYZX+Z^2YX+3Z^4,\\
 W={}&YZ^4Y+ZYZ^3Y+Z^2Y^2Z^2+2Z^2YZ^2Y
      -Z^4Y^2+Z^5X\\
  &-ZYZ^2-2Z^2YZ-2Z^3Y+Z^2 .
\end{aligned}
\]
以下の有限計算を用いる：
(i) $g,U,V,W$の三生成元との交換子は$\mathbb Q$上で零、
(ii) $F_6R_1$での中心核は
$1,g,g^2,U,V,W$が基底となる六次元空間。
これらは命題\ref{prop:quotient-verified-center6}で
有理数演算と法$101$の非零rank証明を組み合わせて証明する。
(iii) 重み$18$以下の$50$中心単項式の像のrankは
法$101$では$49$で、指定の$\Phi$も法$101$で零である。
\emph{特性零における$\Phi=0$は未検証の入力}である。
これを有理数上で厳密に証明できれば、
斉次中心の既知の生成・関係定理と組み合わせて
下記の中心表示を導ける。
さらに斉次楕円Sklyanin代数$S_1$の中心構造定理
\cite{WWYSklyanin}を用いる。
この条件の下で
\[
 Z(R_1)\cong
 \frac{\CC[g,U,V,W]}{(\Phi)},\quad
 \Phi=49U^3-(13g^2+14)U^2+(g^2+1)^2U
 +(11U-g^2-1)VW+V^3-W^3
\]
であり、
$\operatorname{gr}Z(R_1)=Z(S_1)$となる。
\end{proposition}
\begin{proof}[旧原稿の証明と現行の検証範囲]
次数7以下の自由語3280個と関係式倍6015本を
有理数上で消去すると$\dim F_7R_1=120$、
$\dim F_6R_1=84$となる。
交換子写像の法$101$におけるrankは78で、
既に有理数上で中心性・独立性が確認された
上記6元から逆の上界が得られる。
従ってrank78と中心核の基底性は特性零でも確定する
（命題\ref{prop:quotient-verified-center6}）。
斉次中心定理\cite{WWYSklyanin}によれば、
$Z(S_1)$は三次の中心元と三つの六次中心元で生成され、
重み18のただ一つの関係を持つ。
$g,U,V,W$の最高次symbolがそれらを張ることを
命題\ref{prop:quotient-generators-graded}により確認した後、
任意の中心元の最高次項を順に消す次数下降帰納法により
$Z(R_1)$の生成とassociated gradedの等式を得る。

重み18以下の50単項式について、
独立の帰納的語基底アルゴリズムは
法$101$で$\dim F_{18}R_1=1330$、rank49を与える。
$\Phi$を自由語として展開し同じ商へ写した剰余も
法$101$で零と検査した。しかしこれは
\emph{有理数上の零判定の証明ではない}。
特性零における$\Phi=0$を入力として採用すれば、
法$101$のrank49から50単項式の特性零rankは少なくとも49となり、
$\Phi$による非自明な線形関係から高々49。
従って$\Phi$が唯一の重み18の関係式となる。
その最高次symbolは斉次中心の唯一の関係を生成するので、
任意の関係式の最高次項から$\Phi$の倍数を減じる
filtration帰納法により全関係イデアルが$(\Phi)$となる。
\textbf{現在の障害：} 指定した$\Phi$の有理数上の厳密恒等式を
独立に確立する必要がある。法$101$での零判定を
特性零の等式に読み替えてはいけない。
\end{proof}


\begin{proposition}[有限次数中心の独立計算証明]
\label{prop:quotient-verified-center6}
$R_1$の語長filtrationについて
\[
 \dim_{\mathbb C}F_7R_1=120,\qquad
 \dim_{\mathbb C}F_6R_1=84,\qquad
 Z(R_1)\cap F_6R_1=\langle1,g,g^2,U,V,W\rangle_{\mathbb C}.
\]
特に最後の空間の次元は6である。
\end{proposition}
\begin{proof}
次数7以下の自由語は$\sum_{i=0}^7 3^i=3280$個。
三つの定義関係式$r$の全語倍$arb$で
$|a|+|b|\le5$を取ると6015個となる。
これらが定義イデアルの$F_7$部分を生成するので、
有理数上の辞書順疎行列消去で得られる
余次元120は正確な商次元である。
同じ標準剰余から$F_6$の次元84を得る。

付属の\texttt{computations/verify\_center\_degree7.py}
は$\mathbb Q$上で$[g,x_i],[U,x_i],[V,x_i],[W,x_i]$
がすべて零となること、および
$1,g,g^2,U,V,W$が一次独立であることを
元の定義関係式から逐次計算する。
一方、この84次元空間から三生成元との交換子
\[
 F_6R_1\longrightarrow(F_7R_1)^{\oplus3}
\]
の行列を法$101$へ還元して消去するとrank78を得る。
したがって特性零でrankは少なくとも78。
上記の六つの有理数上で独立な中心元により
rankは高々$84-6=78$なので等号である。
計算コードはPython標準ライブラリの有理数演算と
法$101$の整数演算のみで再現可能である。
\end{proof}

\begin{proposition}[全中心の生成と重み18の関係の次元]
\label{prop:quotient-generators-graded}
斉次Sklyanin中心の既知の生成・関係定理
\cite[Proposition~2.31 and \S8.2]{WWYSklyanin}を入力とすると、
\[
 Z(R_1)=\CC[g,U,V,W],\qquad
 \operatorname{gr}Z(R_1)=Z(S_1).
\]
また$\deg g=3$, $\deg U=\deg V=\deg W=6$の
重み18以下の50個の中心単項式が張る空間は、
\emph{特性零でも}次元49である。
したがって重み18以下には1次元の関係空間がある。
ただし、その関係式が上の明示的な$\Phi$と
正確に一致することは、この命題の結論ではない。
\end{proposition}
\begin{proof}
非斉次関係式から定数項を消した斉次Sklyanin代数を
$S_1=\operatorname{gr}R_1$とする。
$g$の三次最高次記号と
$g^2,U,V,W$の六次最高次記号は中心元である。
後の4元が$(S_1)_6$で一次独立であることを
\texttt{computations/verify\_center\_homogeneous\_symbols.py}
の有理数上の計算で確認した。
既知の斉次中心定理では、次数三は一次元、
次数六は四次元で、中心はこれらの元によって生成される。
従って上記の記号は$Z(S_1)$の生成元を与える。

任意の$z\in Z(R_1)$の最高次記号は
$Z(S_1)$の上記生成元の多項式である。
これに対応する$g,U,V,W$の多項式を$z$から引くと、
語長filtrationの次数が下がる。これを有限回繰り返せば
$z$は$g,U,V,W$の多項式となり、
同時に$\operatorname{gr}Z(R_1)=Z(S_1)$を得る。

既知の斉次中心のHilbert級数は
\[
 \operatorname{Hilb}_{Z(S_1)}(t)
 =\frac{1-t^{18}}{(1-t^3)(1-t^6)^3}
\]
である。従って重み18以下の中心部分の次元は
自由可換代数の同範囲の50単項式から、
重み18の一つの関係を引いた49である。
上記の全中心の生成性により、これら50単項式は
$Z(R_1)\cap F_{18}R_1$を張り、
filtered associated gradedの次元保存からrankは49。
残る1次元の関係式の最高次記号は
斉次中心の関係を生成する。しかしその
\emph{非斉次の下位項}は有理数上でまだ同定していない。
\end{proof}

\begin{proposition}[Azumaya locusの同定：既知Sklyanin入力付き]
\label{prop:quotient-azumaya}
命題\ref{prop:quotient-full-center}の中心表示と
$\operatorname{gr}Z(R_1)=Z(S_1)$を仮定し、
$S_1$がPI次数6の三次元AS正則Sklyanin代数で、
その非Azumaya locusが三本の曲線からなること
\cite{WWYSklyanin}、
また$R_1$がfiltered deformationとして
左右大域次元3のprime noetherian整域であることを入力とする。
このとき
\[
 \operatorname{Az}(R_1)
 =\operatorname{Sm}(\operatorname{Spec}Z(R_1)).
\]
特に$R_1$は中心上有限な最大Cohen--Macaulay加群で、
中心上の一般rankは$36$であり、
高さ1ではAzumayaである。
\end{proposition}
\begin{proof}
$Z=Z(R_1)$、$S_1=\operatorname{gr}R_1$とし、
Rees環と中心のRees環を
\[
 \mathcal T=\bigoplus_{d\ge0}F_dR_1\,t^d,\qquad
 \mathcal C=\bigoplus_{d\ge0}(Z\cap F_dR_1)t^d
\]
と書く。元の中心のassociated gradedが
$Z(S_1)$に等しいことから、有限個の斉次中心生成元を
次数を下げつつ持ち上げられる。
従って$\mathcal T$は$\mathcal C$上有限で、
\[
 \mathcal T/(t)=S_1,\quad
 \mathcal C/(t)=Z(S_1),\quad
 \mathcal T[t^{-1}]=R_1[t,t^{-1}],\quad
 \mathcal C[t^{-1}]=Z[t,t^{-1}]
\]
が成り立つ。
$\mathcal C_{(t)}$は一次元noetherian局所整域で
極大イデアルが$t$で生成されるためDVR。
$\mathcal T_{(t)}$は有限torsion-freeゆえ自由で、
そのrankは特殊ファイバー$S_1$の中心上の一般rank
$6^2=36$に一致する。
従って$R_1$の$Z$上の一般rankも$36$。

Zariskian filtrationのgrade/GK次元比較
\cite{LevasseurNoncomm}により、
有限生成$R_1$加群$M$について
\[
 j_{R_1}(M)=j_{S_1}(\operatorname{gr}M),\quad
 \operatorname{GKdim}M=\operatorname{GKdim}\operatorname{gr}M.
\]
$S_1$はAuslander-regularかつGK-Cohen--Macaulay
なので$j_{R_1}(M)+\operatorname{GKdim}M=3$。
中心上有限であるからKrull次元とGK次元が一致し、
Brown--Macleodの基準\cite[Theorem~4.8]{BrownMacleod}
によって$R_1$は$Z$-Macaulayである。

$\mathcal C$の$(t)$を含む閉点で$S_1$がAzumayaなら、
$\mathcal T/\mathfrak m\mathcal T\simeq\operatorname{Mat}_6(\CC)$。
Nakayamaと一般rank36から
$\mathcal T_{\mathfrak m}$は$\mathcal C_{\mathfrak m}$上自由rank36で、
剰余体で同型となる
$\mathcal T_{\mathfrak m}\otimes_{\mathcal C_{\mathfrak m}}
 \mathcal T_{\mathfrak m}^{\mathrm{op}}
\to\operatorname{End}_{\mathcal C_{\mathfrak m}}(\mathcal T_{\mathfrak m})$
はNakayamaにより同型。
従って特殊中心のAzumaya点はRees族でもAzumayaである。

$\mathcal T$の非Azumaya locusを$D$とする。
次数からの$\mathbb G_m$作用で$D$は不変。
$D\cap V(t)$は$S_1$の非Azumaya locusの三曲線で次元1。
$t\ne0$と交わる$D$の既約成分の閉包は
$V(t)$にも交わり、次元公式によりその次元は高々2。
$t\ne0$で$D$は
$(\operatorname{Spec}Z\setminus\operatorname{Az}(R_1))
\times\mathbb G_m$なので、
$R_1$の非Azumaya locusは中心3次元において
次元高々1、すなわちheight one Azumaya。
Brown--Macleod\cite[Theorem~3.13]{BrownMacleod}の
prime、中心有限、有限大域次元、$Z$-Macaulay、
height one Azumayaの条件が揃い、
Azumaya locusと中心のsmooth locusが一致する。
\end{proof}

\begin{lemma}[Azumaya環と中心の大域次元]
\label{lem:quotient-azumaya-gldim}
$C$を可換noetherian環、$B$を$C$上定数正rankの
Azumaya代数とする。このとき左右大域次元は
$\operatorname{gldim}B=\operatorname{gldim}C$であり、
無限大の場合も含む。
\end{lemma}
\begin{proof}
separability idempotentにより
任意の$B$加群$M$は$B\otimes_CM$の$B$直和因子。
$B$は$C$上射影的なので
$C$-射影分解を誘導して
$\operatorname{gldim}B\le\operatorname{gldim}C$を得る。
逆向きには中心への正規化traceで$C\hookrightarrow B$が
$C$加群として分裂する。
任意の$C$加群$N$は$B\otimes_CN$の
$C$直和因子となり、$B$-射影加群を制限すれば
$C$-射影的なので逆の不等式も得る。
\end{proof}

\begin{corollary}[可換中心ファイバーからの十分条件と一つの非正則値]
\label{cor:quotient-azumaya-fibers}
命題\ref{prop:quotient-azumaya}の仮定の下で、
$D_\gamma=R_1/(g-\gamma)$と置く。
$\gamma^2\notin\{0,-1,-28/27\}$なら
$\operatorname{gldim}D_\gamma=2$。
一方、$\gamma^2=-28/27$では
$\operatorname{gldim}D_\gamma=\infty$。
$\gamma=0$での正則性は別途
命題\ref{prop:local-model}の非可換$A_1$局所模型を用いる。
\end{corollary}
\begin{proof}
可換中心ファイバーの必要十分判別式
$H(H+1)(27H+28)$は
\href{https://github.com/uedakazushi/nc-cubic2/blob/main/papers/central_fiber_discriminant.tex}
{別の現行TeX}に証明がある。
$H=\gamma^2$が指定三値以外なら中心曲面は滑らか。
中心の正則元で割った商が正則局所環なら
元の中心局所環も正則なので、
各点で$R_1$はAzumayaであり、
$D_\gamma$は滑らかな二次元中心上Azumaya。
補題\ref{lem:quotient-azumaya-gldim}から結論する。

$\gamma^2=-28/27$では中心曲面の
$(U,V,W)=(1/189,0,0)$は通常二重点だが、
全中心の方程式の$\partial_g$は
$2\gamma\,\partial_HF_H\ne0$なので全中心は滑らか。
従ってこの点で$D_\gamma$は
特異な可換局所曲面上Azumayaとなり、
補題\ref{lem:quotient-azumaya-gldim}から無限大域次元。
\end{proof}
